\section{Anforderungsdefinition}

Aus der in der Einleitung beschriebenen Zielsetzung und den angestrebten Kernfunktionen wurden funktionale sowie nicht-funktionale Anforderungen an das System abgeleitet.

\subsubsection{Funktionale Anforderungen}

Das System muss über ein sprachbasiertes Ein- und Ausgabesystem verfügen, das gesprochene Anfragen erfassen und verarbeitete Antworten als Sprache ausgeben kann.
Eine Wakeword-Erkennung ist erforderlich, um das System gezielt aktivieren zu können, ohne dass eine dauerhafte Verarbeitung aller Audioeingaben stattfindet.
Darüber hinaus muss das System eine API-basierte Erweiterbarkeit bieten, über die zusätzliche Funktionsmodule in das bestehende System integriert werden können.
Auf dem integrierten Display soll ein visueller Avatar dargestellt werden, der den aktuellen Systemzustand widerspiegelt und die Interaktion durch eine grafische Benutzeroberfläche ergänzt.

\subsubsection{Nicht-funktionale Anforderungen}

Die Bedienung des Systems muss intuitiv gestaltet sein, sodass keine technischen Vorkenntnisse für die Nutzung erforderlich sind.
Der Systemaufbau soll modular konzipiert werden, um einzelne Komponenten unabhängig voneinander austauschen oder erweitern zu können.
Die verwendete Hardware muss kompakt dimensioniert sein, damit das Gesamtsystem als Schreibtischgerät eingesetzt werden kann.
Die Integration der einzelnen Hard- und Softwarekomponenten muss robust umgesetzt werden, um einen zuverlässigen Dauerbetrieb zu gewährleisten.
Das System soll echtzeitnahe Reaktionszeiten aufweisen, damit die sprachliche Interaktion als natürlich und flüssig wahrgenommen wird.
Die gesamte Interaktion, einschließlich der sprachlichen Ausgabe, muss verständlich und nachvollziehbar gestaltet sein.
